% -------------------------------------------------
% Academic CV template information
%
% Author: Iacopo Benesperi
% Project page: https://github.com/iacchi/academic_cv
%
% WARNING:
% If you are using this template within Overleaf,
% at the moment this document needs to be compiled
% with either xelatex or lualatex, because the TU
% format for fonts is required by the package
% academicons. The package author says usage with
% pdflatex (T1 fonts) has been fixed already (in
% version 1.9.6), but Overleaf is still running
% version 1.9.1 of the package in their system.
%
% The layout of Academic CV is inspired by My CV,
% work of Pietro Di Leo. The code however is
% almost completely different, offering a more
% streamlined and easier information filling
% compared to the original.
% -------------------------------------------------

% The option SidebarWidth handles the definition
% of the first page sidebar's width, so that it can
% be easily changed to adapt to longer text or a
% different page format. The option has a default
% value of 8.6cm so, unless you want to change the
% default width, you can safely remove the option
% from the declaration. All other declared options
% will be passed directly to the "article" class. 
\documentclass[a4paper,11pt,SidebarWidth=8.6cm]{academic_cv}

% -------------------------------------------------
% Add here extra packages that you may need.
% The academic_cv class already loads the following:
% geometry, libertine, libertinust1math, fontenc,
% fontawesome7, academicons, twemojis, hyperref,
% xcolor, nopageno, graphicx, enumitem, calc, flowfram,
% array, tikz, titlesec, refcount, tabularx.
% -------------------------------------------------



% -------------------------------------------------
% Set colours for header and sidebar in hex code
% -------------------------------------------------

\definecolor{HeaderColor}{HTML}{404040}
\definecolor{HeaderTextColor}{HTML}{FFFFFF}
\definecolor{SidebarColor}{HTML}{E6E6E6}

% -------------------------------------------------
% Header information
% -------------------------------------------------

\Name{John}
\Surname{Doe}
\JobDescription{\LaTeX{} coder extraordinaire}
\RightText{Curriculum Vitae}

% -------------------------------------------------
% Sidebar information
% -------------------------------------------------

% Profile picture, leave empty if you don't want one
\Picture{profile-picture}

% \ContactLangSpacing: Adjust spacing (in mm)
% for the elements of the Contacts, Languages,
% \ConferenceSpacing: Adjust spacing (in pt)
% for the elements of the Conference attendance.
% section.
\ContactLangSpacing{1.2mm}
\ConferenceSpacing{1.2}

\ContactTitle{Contacts}
\ContactDetails{
  % The syntax is: icon in first parenthesis,
  % details in the second parenthesis. Icons
  % using fontawesome or academicons.
  % Date of birth
  \contactlangentry{\faPerson}{Date of birth: XX/XX/XXXX}
  % Primary email
  \contactlangentry{\faEnvelope}{email@address.org}
  % Secondary email
  \contactlangentry{\faEnvelope[regular]}{email@address.it}
  % Phone number
  \contactlangentry{\faPhone}{+XX X XXX XXX XXX}
  % Address
  \contactlangentry{\faLocationDot}{XXX Name Street, City, Post code, Country}
  % ORCID profile
  \contactlangentry{\aiOrcid}{\href{https://orcid.org/0000-0000-0000-0000}{ORCID: 0000-0000-0000-0000}}
  % ResearcherID profile
  \contactlangentry{\aiResearcherID}{\href{https://www.webofscience.com/wos/author/record/P-0000-0000}{ResearcherID: P-0000-0000}}
  % Scopus profile
  \contactlangentry{\aiScopus}{\href{https://www.scopus.com/authid/detail.uri?authorId=00000000000}{Scopus: 00000000000}}
}

\LanguagesTitle{Languages}
\LanguagesDetails{
  % The syntax is: icon in first parenthesis,
  % details in the second parenthesis. Icons
  % using twemojis. The \Large command makes
  % the icons big enough.
  \contactlangentry{{\Large\texttwemoji{it}}}{Language1 - Native language}
  \contactlangentry{{\Large\texttwemoji{uk}}}{Language2 - Professional knowledge (C2)}
  \contactlangentry{{\Large\texttwemoji{fr}}}{Language3 - Basic knowledge (A1)}
}

\ProfSkillsTitle{Professional skills}
\ProfSkillsDetails{
  \skilltag{Skill1}
  \skilltag{Skill2}
  \skilltag{Skill3}
  \skilltag{Skill4}
  \skilltag{Skill5}
  \skilltag{Skill6}
  \skilltag{Skill7}
  \skilltag{Skill8}
}

\SoftSkillsTitle{Soft skills and strengths}
\SoftSkillsDetails{
  \skilltag{Skill1}
  \skilltag{Skill2}
  \skilltag{Skill3}
  \skilltag{Skill4}
  \skilltag{Skill5}
  \skilltag{Skill6}
  \skilltag{Skill7}
  \skilltag{Skill8}
}

\DigitalSkillsTitle{Digital competence}
\DigitalSkillsDetails{
  \skilltag{Skill1}
  \skilltag{Skill2}
  \skilltag{Skill3}
  \skilltag{Skill4}
  \skilltag{Skill5}
  \skilltag{Skill6}
  \skilltag{Skill7}
  \skilltag{Skill8}
}

\OtherInterestsTitle{Other interests}
\OtherInterestsDetails{
  \myinterest{\faMusic{} Music}
  \myinterest{\faDribbble{} Sport}
  \myinterest{\faBook{} Reading}
  \myinterest{\faTrain{} Travels}
  \myinterest{\faCamera{} Photography}
}

\ReferencesTitle{References}
\ReferencesDetails{
  % The syntax is: icon in first parenthesis,
  % details in the second parenthesis. Icons
  % using fontawesome.
  \contactlangentry{\faPerson}{Prof. John Doe -- Venus University \\ john.doe@venus.uni}
  \contactlangentry{\faPersonDress}{Prof. Jane Doe -- University of Mars \\ jane.doe@mars.edu}
}

% -------------------------------------------------
% Main document
% -------------------------------------------------

\begin{document}

\section*{\faGraduationCap{}~~~EDUCATION}

\EduWorkEntry
  {X/XX/XXXX – X/XX/XXXX}
  {PhD in Something}
  {Somewhere, Venus}
  {Venus University}
  {\textit{Thesis title:} Something \\
  \textit{Thesis supervisors:} Prof. John Doe, Prof. Mario Doe \\
  Some explanation.}

\EduWorkEntry
  {X/XX/XXXX – X/XX/XXXX}
  {Master's degree in \\Something long}
  {Somewhere, Mars}
  {University of Mars}
  {\textit{Thesis title:} Something \\
  \textit{Thesis supervisors:} Prof. Jane Doe \\
  Some explanation.\\
  \textbf{Degree mark: XX/XX}}

\EduWorkEntry
  {X/XX/XXXX – X/XX/XXXX}
  {Bachelor's degree in \\Something else}
  {Somewhere, Earth}
  {University of Earth}
  {\textit{Thesis title:} Something \\
  \textit{Thesis supervisors:} Prof. Mario Kart \\
  Some explanation.\\
  \textbf{Degree mark: XX/XX}}

\section*{\faBriefcase{}~~~WORK EXPERIENCE}

\EduWorkEntry
  {XX/XXXX – current}
  {Current title}
  {Somewhere, Earth}
  {University of Earth}
  {Some description.}

\EduWorkEntry
  {XX/XXXX – XX/XXXX}
  {Previous title}
  {Somewhere, Moon}
  {Moon University}
  {Some description.}

\EduWorkEntry
  {XX/XXXX – XX/XXXX}
  {Previous title}
  {Somewhere, Moon}
  {Moon University}
  {Some description.}

\EduWorkEntry
  {XX/XXXX – XX/XXXX}
  {Previous title}
  {Somewhere, Moon}
  {Moon University}
  {Some description.}

\EduWorkEntry
  {XX/XXXX – XX/XXXX}
  {Previous title}
  {Somewhere, Moon}
  {Moon University}
  {Some description.}

\EduWorkEntry
  {XX/XXXX – XX/XXXX}
  {Previous title}
  {Somewhere, Moon}
  {Moon University}
  {Some description.}

\EduWorkEntry
  {XX/XXXX – XX/XXXX}
  {Previous title}
  {Somewhere, Moon}
  {Moon University}
  {Some description.}

\EduWorkEntry
  {XX/XXXX – XX/XXXX}
  {Previous title}
  {Somewhere, Moon}
  {Moon University}
  {Some description.}

\EduWorkEntry
  {XX/XXXX – XX/XXXX}
  {Previous title}
  {Somewhere, Moon}
  {Moon University}
  {Some description.}

% A \newpage is required to separate the
% last \EduWorkEntry of the first page and
% the first \EduWorkEntry of the second
% page. This is to avoid weird margins.
% It is also the only time you have to
% declare a \newpage.
\newpage

\EduWorkEntry
  {XX/XXXX – XX/XXXX}
  {Previous title}
  {Somewhere, Moon}
  {Moon University}
  {Some description.}

\section*{\faBuildingColumns{}~~~TEACHING EXPERIENCE}

\EduWorkEntry
  {XX/XXXX - XX/XXXX}
  {Teacher}
  {Somewhere, Earth}
  {University of Earth}
  {I taught something.}

\section*{\faTrophy{}~~~FUNDING AND AWARDS}

\begin{itemize}[leftmargin=*,itemsep=0pt,parsep=0pt]
  \item{I got some big money.} 
  \item{I got some medium money.} 
  \item{I got some small money.}
\end{itemize}

\section*{\faBookOpen{}~~~PUBLICATIONS}

Scientific parameters: h-index: XX; total citations: XXXX. (source: Web of Science/Scopus)
\begin{enumerate}[leftmargin=*,itemsep=0pt,parsep=0pt,topsep=0pt]
  \item \publication{2025}{\textbf{J. Doe}, M. Kart, et al.}{Some title}{Journal of nothing}{10}{600-700}{10.1000/j.not.314864}{JIF: 18.9. Citations: XX}
\end{enumerate}

\section*{\faPeopleGroup{}~~~CONFERENCE ATTENDANCE}

% This text goes into a tabularx environment. The
% syntax is: type of contribution (Poster, Talk,
% Attendance, Invited, Keynote) in first column,
% details in the second column.
\Conferences{
  \textbf{Keynote} & \textbf{J. Doe}, M. Kart, et al. ``Some fantastic title'', XXX Conference 2026, Somewhere, Earth, XX/XXXX.\\
  \textbf{Invited} & \textbf{J. Doe}, M. Kart, et al. ``Another great title'', Moon conference 2025, Somewhere, Moon, XX/XXXX.\\
  \textbf{Talk} & \textbf{J. Doe}, M. Kart, et al. ``My first conference talk'', 2024 Mars Meeting, Somewhere, Mars, XX/XXXX.\\
  \textbf{Poster} & \textbf{J. Doe}, M. Kart, et al. ``A decent title'', MercuryCon 2023, Somewhere, Mercury, XX/XXXX.\\
  \textbf{Attendance} & Pluto Conference 2023, Very far, Pluto, XX/XXXX.
}

\section*{\faCertificate{}~~~SCIENTIFIC LEADERSHIP}

\subsection*{Organisation of international conferences}

\begin{itemize}[leftmargin=*,itemsep=0pt,parsep=0pt,topsep=0pt]
  \item{I organised a conference!}
\end{itemize}

\subsection*{Supervision and mentoring activities}

\begin{itemize}[leftmargin=*,itemsep=0pt,parsep=0pt,topsep=0pt]
  \item{I supervised someone!}
\end{itemize}

\subsection*{Outreach}

\begin{itemize}[leftmargin=*,itemsep=0pt,parsep=0pt,topsep=0pt]
  \item{\textbf{Event title} XXXX \\ I was part of the organising team of Event title.}
\end{itemize}

\subsection*{Academic peer review}

\begin{itemize}[leftmargin=*,itemsep=0pt,parsep=0pt,topsep=0pt]
  \item{I reviewed manuscripts for the Journal of nothing.}
\end{itemize}

\subsection*{Other activities}

\begin{itemize}[leftmargin=*,itemsep=0pt,parsep=0pt,topsep=0pt]
  \item{I did something else as well.}
\end{itemize}

\end{document}